\scanpage{15}
Ceci posé, on \unsure{suppose} $W$ construit en recollant $P$ et $Q$ suivant leurs bords \deleted{\unsure{communs}} via $\varphi$.

\deleted{\unsure{On peut}}

\deleted{\unsure{On aura}}

L'hypothèse que $\partial P$ est un fibré trivial sur $S$ implique que l'inclusion \deleted{\unsure{un}}
\diagram{\begin{tikzcd}
W-Q=P-\partial P \arrow[r,hook] & P
\end{tikzcd}}
est un isomorphisme d'homotopie,

\deleted{\unsure{En fait}}

et de même pour
\diagram{\begin{tikzcd}
W-P=Q-\partial Q \arrow[r,hook] & Q & \text{''\qquad ''\qquad ''}.
\end{tikzcd}}

\deleted{\unsure{On a donc}}

Considérons le diagramme
\diagram{\begin{tikzcd}[row sep=large,column sep=large]
& & T=\partial M & t \arrow[l,very thick] \\
& W & Q \arrow[l,hook] \arrow[u,"\unsure{\simeq|2}"] & N=D_{\unsure{1}}\sim e \arrow[l,hook] \arrow[u,hook,very thick] \\
\partial N=S^1 & P \arrow[l] \arrow[u,hook,very thick] & {\begin{gathered}S^1\times T=\partial N\times\partial M\deleted{\unsure{\partial}}\\=\partial P=\partial Q=P\cap Q\end{gathered}} \arrow[l,hook] \arrow[u,hook] & \partial N=\unsure{S^1} \arrow[l,hook] \arrow[u,hook] \\
s \arrow[u] & M \arrow[l] \arrow[u,hook] & \partial M=T \arrow[l,hook] \arrow[u,hook] & (s,t) \arrow[l,hook] \arrow[u,hook]
\end{tikzcd}}
\mspage{14}
